\documentclass[11pt,letterpaper]{article}
\usepackage[T1]{fontenc}
\usepackage{lmodern,microtype,amsmath,amssymb,amsthm,booktabs,tabularx}
\usepackage[margin=1.05in]{geometry}
\usepackage{xcolor}
\definecolor{inkblue}{RGB}{22,63,97}
\usepackage[pdfusetitle,colorlinks=true,linkcolor=inkblue,urlcolor=inkblue,citecolor=inkblue,
pdftitle={Unbounded minimum Scott ranks for Boolean algebras}]{hyperref}
\newtheorem{theorem}{Theorem}[section]
\newtheorem{lemma}[theorem]{Lemma}
\newtheorem{proposition}[theorem]{Proposition}
\newtheorem{corollary}[theorem]{Corollary}
\theoremstyle{definition}\newtheorem{definition}[theorem]{Definition}
\newcommand{\N}{\mathbb N}
\newcommand{\Baire}{\mathcal N}
\newcommand{\BA}{\mathsf{BA}}
\newcommand{\Clop}{\operatorname{Clop}}
\newcommand{\SR}{\operatorname{SR}}
\newcommand{\Hit}{\operatorname{Hit}}
\newcommand{\rk}{\operatorname{rk}}
\newcommand{\id}{\mathrm{id}}
\newcommand{\card}[1]{\lvert #1\rvert}
\newcommand{\res}{\mathbin{\upharpoonright}}
\newcommand{\code}[2]{\href{https://github.com/SerafimBatzoglou/boolean-scott-spectral-gaps/blob/e1b9ab53208242a9da1efc2336e61db9c1a5566c/BooleanScott/#1.lean}{#2}}
\setcounter{tocdepth}{1}
\title{Unbounded minimum Scott ranks\\for Boolean algebras}
\author{GPT 6.2 Astra (Codex) with Serafim Batzoglou}
\date{Research manuscript\\November 26, 2026}
\begin{document}
\maketitle
\begin{abstract}
For every countable ordinal \(b\) we construct, by an existence argument, a satisfiable
pure Boolean \(\Pi_{\omega+4}\) sentence none of whose countable models admits a
\(\Pi_{b+1}\) Scott sentence. This answers the existence part of Question~1.9
of Gonzalez and Harrison-Trainor negatively at a fixed input complexity.
We build a countable Boolean algebra \(B\) and associate a sentence to each tree
on \(\mathbb N\). Every countable model \(A\) of that sentence satisfies
\(A\cong A\times B\) exactly when the tree has an infinite branch. The
criterion follows from finite counts in clopen sets on the Stone side and a
homeomorphism constructed along a branch. Projecting a Borel set of operation
tables then gives the lower bound simultaneously for all models of one sentence.
\end{abstract}
\tableofcontents
\clearpage

\section{The minimum of a Scott spectrum}\label{sec:intro}
Scott's isomorphism theorem gives every countable structure \(A\) a sentence of
\(L_{\omega_1\omega}\) whose countable models are precisely its isomorphic
copies~\cite{scott}. The complexity of the simplest such description measures
how difficult it is to distinguish \(A\) from the other countable structures.
A different question arises when we specify a \emph{class} of structures.
How difficult must it be to describe its simplest member?

We use the parameterless Scott rank of Montalb\'an~\cite{montalban} in its
Scott-sentence formulation. Thus \(\SR(A)\) is the least ordinal \(b\) for
which \(A\) has a \(\Pi_{b+1}\) Scott sentence. The Scott spectrum of a sentence
is the set of these ranks for its countable models. Without a restriction on
the class of structures, Harrison-Trainor showed that a satisfiable
\(\Pi_2\) sentence can have arbitrarily high minimum rank~\cite{ht}.
For linear orders, Gonzalez and Harrison-Trainor proved that every satisfiable
\(\Pi_\alpha\) input extending the linear-order axioms has a model with a
\(\Pi_{\alpha+4}\) Scott sentence~\cite[Theorem~1.5]{ght}.
They ask whether any bound depending on \(\alpha\) exists for Boolean
algebras~\cite[Question~1.9]{ght}. The question concerns the minimum of the
spectrum. Exhibiting one complicated Boolean algebra cannot answer it.

\begin{theorem}\label{thm:main}
For each countable ordinal \(b\) there is a sentence \(\varphi_b\) in the
signature \(0,1,\neg,\wedge,\vee\), with equality, such that
\begin{enumerate}
\item \(\varphi_b\) is \(\Pi_{\omega+4}\), and every structure satisfying it
is a Boolean algebra with \(0\ne1\);
\item \(\varphi_b\) has a countable Boolean model;
\item for every countable \(A\models\varphi_b\), no \(\Pi_{b+1}\) sentence
isolates \(A\) up to isomorphism among all countable structures in this signature.
\end{enumerate}
\end{theorem}
In particular, the set of minima of Scott spectra of these satisfiable inputs
is unbounded in \(\omega_1\). The witness sentence depends on \(b\). Our
argument gives an upper bound on its syntactic complexity and leaves the first
possible level of failure undetermined.

\subsection*{Why use a product?}
If a structure agrees with a nonisomorphic structure on every
\(\Pi_{b+1}\) sentence, its rank exceeds \(b\). We arrange that a comparison
structure is always available: it is the product \(A\times B\) with a
\emph{fixed} Boolean algebra \(B\). The product corresponds under Stone
duality to adjoining a fixed disjoint summand \(R\). Absorption of summands
is a delicate invariant of countable Boolean algebras. For linear orders,
the existence in \(M\) of initial segments of type \(L\cdot n\) for all
finite \(n\) forces \(L+M\cong M\). Gonzalez and Harrison-Trainor
explain how Ketonen's work on Boolean algebras~\cite{ketonen} precludes
the analogous inference there~\cite[\S1]{ght}. We use this flexibility
to encode the existence of an infinite branch.

The proof has two distinct quantifiers to control. First we construct a
sentence \(\Psi_T\) for each tree \(T\), with an explicit countable witness,
and prove
\begin{equation}\label{eq:criterion}
 \text{for every countable Boolean } A\models\Psi_T,\qquad
 A\cong A\times B\quad\Longleftrightarrow\quad T\text{ has a branch}.
\end{equation}
The sentence prescribes small clopen neighborhoods of certain points of the
Stone space. It allows additional points elsewhere. Sections~\ref{sec:guard}
and~\ref{sec:absorb} explain why such points cannot invalidate either
direction of~\eqref{eq:criterion}. Second, for a given \(b\), we must find
one branchless \(T\) for which \emph{every} model agrees with its product
through level \(b+1\). The existence of even one model that fails this
comparison is an analytic condition on \(T\). The branchless parameters form
a nonanalytic set. Section~\ref{sec:diagonal} uses this discrepancy to choose \(T\).

\section{Logic and topology used in the construction}\label{sec:tools}
Write \(\N=\omega=\{0,1,\ldots\}\), and take all ordinal bounds below the first
uncountable ordinal \(\omega_1\). Countable includes finite.
The language \(L_{\omega_1\omega}\) allows countable conjunctions and
disjunctions and finite strings of quantifiers, with finitely many free
variables per formula. We use the usual noneffective infinitary hierarchy.
At level zero take finitary quantifier-free formulas. For \(\alpha>0\),
the two forms are
\[
 \Sigma_\alpha:\ \bigvee_i\exists\bar u_i\,\theta_i,\quad
       \theta_i\in\Pi_{\beta_i};
 \qquad
 \Pi_\alpha:\ \bigwedge_i\forall\bar u_i\,\eta_i,\quad
       \eta_i\in\Sigma_{\beta_i},\qquad \beta_i<\alpha.
\]
The index sets are countable, the blocks finite (possibly empty), and a
quantifier-free formula is available at every level. Negation exchanges the
two forms. At a positive level one may flatten countable conjunctions of
\(\Pi\)-formulas, or countable disjunctions of \(\Sigma\)-formulas, at that
same level. Finite Boolean combinations of finite-level formulas have finite
level. A countable collection of tests, each at some finite level, can be
conjoined at \(\Pi_\omega\) by first putting each at a finite \(\Sigma\)
level. A \emph{finite-level} test can itself be infinitary.

A Boolean algebra is a bounded distributive lattice in which every element
\(a\) has a complement \(\neg a\), characterized by \(a\wedge\neg a=0\)
and \(a\vee\neg a=1\). The order is \(a\le b\) iff \(a\wedge b=a\).
A filter is an upward-closed set containing \(1\) and closed under finite
meets; an ultrafilter is a maximal filter omitting \(0\). It picks exactly
one of \(a,\neg a\) for each \(a\).

A \emph{Stone space} is a compact Hausdorff space with a basis of sets
that are both open and closed (clopens). Its clopens form a Boolean algebra
\(\Clop(X)\) under the set operations. Conversely the ultrafilters of a
Boolean algebra \(A\), with basic clopens \(\widehat a=\{p:a\in p\}\), form
its Stone space. The map \(a\mapsto\widehat a\) is an isomorphism:
ultrafilters separate distinct elements, and compactness turns a covering of
a clopen by basic clopens into a finite union. Stone duality~\cite{stone}
identifies isomorphisms of these algebras with homeomorphisms of their
spaces. In particular,
\[
 \Clop(X\sqcup Y)\cong\Clop(X)\times\Clop(Y),\qquad
 \Clop(C)\cong[0,C]\subseteq\Clop(X)\quad(C\text{ clopen}).
\]
The latter interval has top \(C\) and complement \(x\mapsto C\setminus x\).
Any formula about the interval can be translated to one about the ambient
algebra with the additional parameter \(C\), at the same level. Replace
variables by \(x\wedge C\), top by \(C\), and complements by relative
complements. Every interval assignment lifts to an ambient assignment,
so induction through either kind of quantifier block proves preservation
of truth. This also covers \(C=\varnothing\).
We use Stone spaces with countably many clopens throughout the construction.
They have countable bases, although their sets of points can be uncountable.

For compact Stone spaces \(X_i\), put
\[
 \mathcal K(X_i:i\in I)=\{\infty\}\ \cup\ \coprod_{i\in I}X_i .
\]
Give the disjoint union its sum topology and \(\infty\) the neighborhoods
of the one-point compactification. The components are clopen. Neighborhoods
of \(\infty\) contain every whole component outside some finite set:
a compact subset of a sum meets only finitely many components. Conversely
the complements of finite unions of whole components form a neighborhood
basis there. Hence a clopen of \(\mathcal K\) is specified by its membership
bit at \(\infty\), and by a clopen in each component agreeing with that
bit (whole or empty) outside finitely many indices. These data also prove
compactness, zero-dimensionality and countability of the clopen algebra
when \(I\) and each \(\Clop(X_i)\) are countable. We call the added point
the \emph{apex}. For a set \(P\subseteq X\), its derivative \(P'\) is
the set of all points every neighborhood of which meets \(P\) in a point
other than itself. Set \(P^{(0)}=P\), \(P^{(k+1)}=(P^{(k)})'\).
At an apex, belonging to \(P'\) is exactly the condition that \(P\)
meets infinitely many components.

\subsection*{Translating incidence into Boolean formulas}
Our point sets will be intrinsic: restricting the ambient space to an open
subspace restricts the point set by intersection. Isolated points, closure,
derivative, and the set operations used below all have this locality.
For such a set \(P\), write \(\Hit_P(C)\) for \(C\cap P\ne\varnothing\).
We need tests using quantified \emph{clopens}; the language has no point
variables.

\begin{lemma}\label{lem:hit}
Suppose \(\Hit_P\) has a finite-level Boolean test on clopens.
For every \(h<\omega\), the condition \(\card{C\cap P}\ge h\) has a
finite-level test whose level is bounded independently of \(h\).
Incidence with \(\overline P\) and with \(P'\) also has a finite-level
test. Given finite-level incidence tests for \(P,Q\), one can test
\(P\setminus Q\) if \(Q\) is closed, and \(P\cap Q\) if both are closed.
\end{lemma}
\begin{proof}
For \(h>0\), ask for \(h\) pairwise disjoint clopens below \(C\), each
meeting \(P\). Distinct points in \(C\) can be separated by such clopens:
take a finite clopen partition distinguishing the chosen points. The
converse uses disjointness; \(h=0\) is a tautology. There is one finite
existential block regardless of \(h\), and the finitely many incidence
tests in its matrix can all be put at one common positive \(\Pi\) level.
Thus the bound is uniform. Incidence with \(\overline P\) equals incidence
with \(P\) since \(C\) is open. A clopen meets \(P'\) exactly when it
contains infinitely many points of \(P\). One direction follows by separating
finitely many points in a Hausdorff space; the other uses compactness of
\(C\). Conjoin all the threshold tests at one further finite level.

To test \(P\setminus Q\), ask for a subclopen meeting \(P\) and missing
\(Q\). Closedness gives such a neighborhood of any witness. Finally
\(C\cap P\cap Q\) is empty exactly when some clopen \(U\) contains
\(C\cap P\) and misses \(C\cap Q\). For disjoint compact closed sets this
follows by covering the first with finitely many clopens disjoint from
the second. The condition can be written
\(\neg\Hit_P(C\setminus U)\wedge\neg\Hit_Q(C\cap U)\).
All quantification is over elements of \(\Clop(X)\).
\end{proof}

\subsection*{From tuple orbits to a Scott sentence}
An orbit of a finite tuple consists of its images under all automorphisms.
For \(\bar a=(a_1,\dots,a_m)\) of clopens, its \emph{sign cells} are the
\(2^m\) intersections obtained by taking either \(a_i\) or its complement
at each position. Empty cells are allowed. The nonempty cells partition
\(X\), and a homeomorphism on each pair of corresponding cells glues to a
homeomorphism taking \(\bar a\) to \(\bar b\). Each cell is given by a
Boolean term, so every finite-level test on one cell gives a finite-level
test on the tuple by substitution. For \(m=0\) the only cell is \(X\).

\begin{lemma}\label{lem:orbits}
Let \(A\) be a countable Boolean algebra. Suppose that for each finite
tuple \(\bar a\) of \(A\) there is a \(\Pi_\lambda\) formula \(D_{\bar a}\)
defining its orbit \emph{in \(A\)}. Then \(A\) has a \(\Pi_{\lambda+2}\)
Scott sentence among countable Boolean algebras. In particular, if agreement
on all finite-level formulas implies that two tuples in \(A\) are
automorphic, \(A\) has a \(\Pi_{\omega+2}\) Scott sentence.
\end{lemma}
\begin{proof}
Require \(D_\varnothing\). For every \(\bar a\), require that any tuple
\(\bar x\) satisfying \(D_{\bar a}\) have the same order relations between
its entries as \(\bar a\). Add the following two extension clauses:
\[
\begin{aligned}
 &\forall\bar x\bigl(D_{\bar a}(\bar x)\ \longrightarrow\
        \exists y\,D_{\bar a a}(\bar x,y)\bigr) &&(a\in A),\\
 &\forall\bar x\,y\bigl(D_{\bar a}(\bar x)\ \longrightarrow\
        \bigvee_{a\in A} D_{\bar a a}(\bar x,y)\bigr).
\end{aligned}
\]
The bodies of the universal clauses are \(\Sigma_{\lambda+1}\):
negate the antecedent and take a disjunction. Countability of the tuples
and one outer universal block give \(\Pi_{\lambda+2}\). In \(A\) an
automorphism witnessing \(D_{\bar a}\), or its inverse, supplies the
extensions. In a countable model, enumerate both algebras and extend a
matched pair of tuples alternately to include each next element. The order
clauses preserve and reflect equality and order, so their union is a
bijective order isomorphism, hence a Boolean isomorphism.

For the last assertion fix \(\bar a\). For each tuple of the same length
outside its orbit choose a distinguishing finite-level formula true of
\(\bar a\), negating it if necessary. Move each test to a finite
\(\Sigma\) level and conjoin. This is a \(\Pi_\omega\) definition of the
orbit inside \(A\), also when the list of exclusions is empty.
\end{proof}
\section{Points of distinguishable finite degrees}\label{sec:packets}
We need infinitely many distinguishable compact spaces, each recognizable
at some finite logical level. We construct them from finite-height scattered
spaces and Cantor space \(K=\{0,1\}^{\N}\) with the product topology. A cylinder
in \(K\) prescribes finitely many initial bits. Write \(E\cong F\) for homeomorphism
and say that \(E\) \emph{absorbs} \(F\) if \(E\sqcup F\cong E\).
A clopen summand of \(E\) that absorbs \(F\) makes \(E\) absorb \(F\)
as well, by using the identity on its complement.

\begin{definition}
Set \(T_0=\{*\}\), \(T_{d+1}=\mathcal K(T_d:i\in\N)\), and define the
\emph{packet of degree \(d\)} by
\[
 G_d=\mathcal K(K\sqcup T_d:i\in\N).
\]
In any Stone space \(X\) let \(D=\overline{\operatorname{Iso}(X)}\), the
closure of its isolated points. Define the degree-\(d\) marker set by
\begin{equation}\label{eq:marker}
 M_d(X)=\bigl(D\cap\overline{X\setminus D}\cap D^{(d+1)}\bigr)
                         \setminus D^{(d+2)}.
\end{equation}
\end{definition}

\begin{lemma}\label{lem:packets}
The set \(M_j(G_d)\) is the singleton consisting of the apex for \(j=d\)
and is empty for \(j\ne d\). Each \(M_j\) is local and has a finite-level
incidence test.
\end{lemma}
\begin{proof}
Isolated points are dense in \(T_d\). Also \(T_d^{(d)}\) is a singleton
and \(T_d^{(d+1)}\) is empty, where the derivative starts from the entire
space. Both assertions follow by induction: each component restriction is
the derivative in \(T_{d-1}\), and the apex survives the next derivative
precisely when infinitely many components still contribute. In \(G_d\)
the set \(D\) consists of the tower components and the apex, and its
complement consists of the Cantor components. Thus \(D\cap\overline{X\setminus
D}\) consists of the apex alone. In fact \(D^{(d+1)}\) is the singleton
consisting of the apex, and \(D^{(d+2)}\) is empty by the same component calculation. This gives
the claimed marker.

A clopen contains an isolated point exactly when it lies above an atom,
a nonzero Boolean element with no nonzero proper part. This gives a
finite-level test for \(\operatorname{Iso}(X)\) and for its closure.
Apply Lemma~\ref{lem:hit} to the successive closed derivatives, to the
complement of \(D\), to its closure and to the closed intersections in
\eqref{eq:marker}, followed by the final difference. All these operations
commute with restriction to open sets. The bound may depend on \(j\).
\end{proof}

\begin{lemma}\label{lem:fragments}
Every clopen in a \(T_d\) is empty or is a finite positive disjoint sum
of copies of a single \(T_q\), \(q\le d\). A clopen of \(G_d\) containing
its apex is homeomorphic to \(G_d\). A clopen missing its apex consists of
an optional copy of \(K\) and an empty or finite positive sum of copies of
one \(T_q\), \(q\le d\). Every one of these homeomorphism types, including
\(G_d\) itself, can be specified among countable clopen algebras by a
finite-level sentence.
\end{lemma}
\begin{proof}
A repeated compactified sum absorbs an extra copy of its repeated component:
shift the component indices and fix the apex. Thus \(T_{q+1}\) absorbs
\(T_q\), and it absorbs every lower tower by applying the induction
hypothesis inside one of its components. In a finite union of towers the
greatest height absorbs all lower heights. The clopen description of
\(\mathcal K\) now proves the assertion for a tower by induction. An
apex-free clopen is a finite union of lower fragments. An apex-containing
one contains a complete tail homeomorphic to the whole tower, with only
finitely many lower fragments left over. Likewise in \(G_d\) the punctured
sum can be rearranged into one infinite list of Cantor components and
one infinite list of towers of height \(d\). Shifting either list proves
absorption of \(K\) and \(T_d\); a tower component absorbs lower towers.
A nonempty clopen of \(K\) is a finite union of cylinders, each homeomorphic
to \(K\), and any finite positive number of Cantor spaces is again \(K\)
by a finite cylinder partition. These observations give the two clopen
forms in \(G_d\).

Here are the logical bounds. Write \(\#P\) for cardinality truncated to
\(\N\cup\{\infty\}\). Within \(G_d\), a clopen sign cell \(U\) containing
the apex is detected by a nonzero count of \(D^{(d+1)}\). For the other
cells the counts of \(D^{(k)}\), \(0\le k\le d\), determine the tower
part. If it is nonempty, the last nonzero stage is its height \(q\) and
that stage counts precisely its multiplicity. The count in the ordinary
\((d+1)\)-st derivative of the whole space detects the Cantor part by
being nonzero. All counts here are taken in \(U\); locality identifies
them with counts in the corresponding loci of the ambient space. In a
fixed tower the ordinary derivative counts up to its height likewise
determine the clopen type. Within \(K\) only emptiness of each cell
matters.

For each reference tuple record every finite threshold satisfied or failed
by each of these counts on each sign cell. These data define the orbit:
equal data give homeomorphic cells, and the cell maps glue. The formulas
have a \emph{common finite bound} for a fixed ambient \(G_d\), \(T_d\)
or \(K\). Indeed the list of loci is finite, the threshold bound in
Lemma~\ref{lem:hit} is independent of the threshold, and substitution of
a sign term does not change the level. Put every assertion or its negation
at a common finite \(\Sigma\) level one step higher if necessary; a single
countable conjunction gives a finite \(\Pi\) level independent of tuple
length. Lemma~\ref{lem:orbits} gives a finite-level Scott sentence in
each case. The empty space is specified by \(0=1\). Finally, if two
spaces have finite-level Scott sentences, their disjoint sum can be tested
by asking for a clopen and its complement satisfying the two relativized
sentences. Stone duality gives the converse. This proves the assertion
for every finite sum appearing in the classification.
\end{proof}

\section{Rows, vertex charts, and their Scott sentences}\label{sec:charts}
The same family of packets serves two purposes. Odd degrees indicate
which rows are present. Even degrees will measure the change when an
extra row is added. For \(e,n\in\N\) and nonempty \(S\subseteq\N\), set
\begin{align}
 R_{e,n}
   &=\left(\;\coprod_{j\ge n}G_{2j}\ \sqcup\
                 \coprod_{k\in\N}G_{2e+1}\right)^+ ,\label{eq:row}\\
 C_S&=\left(\coprod_{e\in S,\ k\in\N}R_{e,e+k}\right)^+ .\label{eq:chart}
\end{align}
Here \(+\) as a superscript means to adjoin the compactification point.
We call \(e\) the \emph{tag}, \(n\) the \emph{cutoff}, the apex of a row
its \emph{root}, and the apex of a chart its \emph{vertex}. A packet of
degree \(2j\) occupies \emph{column \(j\)}. Write
\[
    A_e(X)=M_{2e+1}(X)'',\qquad P_e(X)=M_{2e+1}(X)'\setminus A_e(X).
\]
These are local and have finite-level incidence tests for each \(e\);
\(A_e\) is closed and \(P_e'\subseteq A_e\).
In a chart \(A_e\) will identify the vertex while \(P_e\) identifies
the row roots of tag \(e\).

\begin{lemma}\label{lem:census}
In \(R_{a,n}\), \(P_e\) consists of the root if \(a=e\) and is empty
otherwise; every \(A_e\) is empty. In \(C_S\), \(A_e\) consists of the
vertex if \(e\in S\) and is empty otherwise; \(P_e\) consists exactly of
the roots of rows of tag \(e\). The points of \(M_{2j}\) are precisely the
apices of the packets in column \(j\). In particular their number in a
chart is finite.
\end{lemma}
\begin{proof}
A row root lies in \(D^{(m)}\) for every finite \(m\): arbitrarily late
even packets contain points in each fixed stage of \(D\), and the apex
test gives the next stage at the root. The chart vertex also lies in every
finite stage, since infinitely many row roots do. Thus neither kind of
point belongs to any \(M_j\). All other points are in clopen packets and
Lemma~\ref{lem:packets} applies. In a row the infinitely repeated odd
marker accumulates exactly at its root; there is at most one marker of
each even degree. In a chart the rows containing odd markers of tag \(e\)
are exactly the infinitely many rows with first index \(e\), if any.
Consequently the first derivative consists of their roots together with
the vertex when this tag occurs, and its next derivative consists of the
vertex alone. In column \(j\) the contributing row indices satisfy
\(e+k\le j\), giving at most \((j+1)(j+2)/2\) points.
\end{proof}

\paragraph{A chart with two tags.}
Take \(S=\{0,2\}\). The entry in each column below counts the rows of
the indicated tag that supply a packet in that column.
\begin{center}
\begin{tabular}{@{}lrrrr@{}}
\toprule
Column \(j\) & \(0\) & \(1\) & \(2\) & \(3\)\\
\midrule
Rows of tag \(0\) with cutoff \(\le j\) & 1 & 2 & 3 & 4\\
Rows of tag \(2\) with cutoff \(\le j\) & 0 & 0 & 1 & 2\\
\midrule
Number of points of \(M_{2j}\) & 1 & 2 & 4 & 6\\
\bottomrule
\end{tabular}
\end{center}
For instance, column \(2\) receives one marker from each of the three
tag-\(0\) rows with cutoffs \(0,1,2\), and one from the tag-\(2\) row
of cutoff \(2\). The infinitely many odd packets in each of these rows
contribute nothing to the even count. Conversely, the even counts are
insensitive to the tag of an added row: adjoining one row of cutoff zero
increases \emph{every} even-column count by one. Tracking where this
extra unit can go will force a branch in Section~\ref{sec:absorb}.

\subsection{Matching rows across a family of columns}
\begin{lemma}\label{lem:match}
Consider two disjoint sums of rows with countable row sets. Suppose there
is a tag-preserving bijection of their row sets. The cutoff of a row may
change under the bijection. If in each column the two numbers of eligible
rows (cutoff at most the column index) are finite and equal, then the sums
are homeomorphic by a map inducing the prescribed bijection on roots.
The map extends to their one-point compactifications.
\end{lemma}
\begin{proof}
Identify the row sets by the specified bijection. In column \(j\) let
\(I_j,J_j\) be the two finite eligible sets. Fix the common elements and
choose any bijection \(I_j\setminus J_j\to J_j\setminus I_j\); equality
of the column sizes permits this. Map the even packets using these
bijections and the identity on their model \(G_{2j}\). In each paired
row map its odd packets bijectively, preserving the repeated degree.
Map each root to its partner. This defines a bijection of the whole sums.

For row \(i\) the maps in all columns beyond the maximum of its two
cutoffs fix \(i\), in \emph{both} directions. Only finitely many packets
therefore depart from or enter this row from a different row. Given a
neighborhood of the target root, discard these finitely many departures
and the preimages of the finitely many exceptional packets in the
neighborhood. What remains contains a neighborhood of the source root.
The inverse has the same property. All other points lie in open packets.
A homeomorphism takes compact sets to compact sets in both directions,
so extending it by the point at infinity is continuous with continuous
inverse.
\end{proof}

\subsection{Recognizing an arbitrary clopen in a chart}
Given a clopen \(U\subseteq C_S\), write \(\epsilon(U)=1\) if it contains
the vertex and \(0\) otherwise. Call a row \emph{retained} if its root
belongs to \(U\), and write \(I_U\) for the set of retained row indices.
A portion of \(U\) can cut through a packet; agreement just on which roots
are present would lose this information. We separate the full tails from
the finitely many pieces where such cuts can occur.

\begin{lemma}\label{lem:core}
For every \(U\) there exists \(J_0\) such that for each \(J\ge J_0\)
there is a clopen \(U^*\subseteq U\) with the following description.
For each retained row \(i=(e,k)\), its part in \(U^*\) is a full row of
tag \(e\) and cutoff \(\max(J,e+k)\) up to homeomorphism. There are no
other rows in \(U^*\). Adjoin the outer vertex exactly when
\(\epsilon(U)=1\). For \(j\ge J\), the number of eligible rows in this
description is precisely \(\#(U\cap M_{2j})\).
The remainder \(U\setminus U^*\) is a finite disjoint sum of full mixed
packets \(G_d\) and a space consisting of an optional Cantor part and
finitely many towers of bounded height.
\end{lemma}
\begin{proof}
Outside finitely many rows \(F\), membership in \(U\) is constant on the
entire row, with value \(\epsilon(U)\). In each row of \(F\) the same
description at the row root gives finitely many exceptional packets.
Choose \(J_0\) after all exceptional even columns in these rows. Choose
an odd starting index after their exceptional odd packets, separately
for each exceptional row; in other rows use odd starting index zero.
In a retained row take its root, its even packets in columns \(j\ge J\)
and all its odd packets from the chosen starting index onwards. This
clopen is a row of the asserted type by reindexing its odd packets.
Take the empty set in the other rows. Outside \(F\) only rows with
\(e+k<J\) could have a changed restriction. There are finitely many of
these, hence the indicated vertex bit makes \(U^*\) clopen in the chart.
Our choices ensure \(U^*\subseteq U\). Above the cut \(J\), membership
of every whole even packet in \(U\) is exactly its row-root bit;
Lemma~\ref{lem:census} gives the column count.

The difference is compact, misses the vertex, and hence meets only
finitely many rows. Its restriction to each is compact and misses the
row root, so meets only finitely many packets there. Apply
Lemma~\ref{lem:fragments} in each of these packets to get the statement
about the remainder.
\end{proof}

\begin{proposition}\label{prop:chart}
For every nonempty \(S\subseteq\N\), the Boolean algebra \(\Clop(C_S)\)
has a \(\Pi_{\omega+2}\) Scott sentence among countable Boolean algebras.
\end{proposition}
\begin{proof}
First compare two clopens \(U,V\) of this chart agreeing on every unary
finite-level test. By Lemmas~\ref{lem:hit} and~\ref{lem:census}, they have
the same vertex bit (test \(A_e\) for some \(e\in S\)), the same extended
number of retained roots of each tag (test \(P_e\)), and the same count
of \(M_d\) for each \(d\). These assertions follow from all finite
thresholds; in the infinite case the sets of row indices are countable.

If \(U\) has no retained root, it misses the vertex, since any vertex
neighborhood contains all but finitely many whole rows. The argument for
the remainder in Lemma~\ref{lem:core} expresses \(U\) as a finite sum
of packet fragments. It has a finite-level Scott test by
Lemma~\ref{lem:fragments} and finite splitting. Relativizing that test to
the clopen gives \(V\cong U\), including \(U=\varnothing\).

If there is at least one retained row, choose a common threshold \(J\)
for both applications of Lemma~\ref{lem:core}. For every tag use the
equality of finite cardinalities or a bijection between the two countably
infinite fibers to pair all retained roots. The two cores have zero
eligible rows below \(J\). At and above \(J\) they have equal finite
column counts by the lemma. Lemma~\ref{lem:match} matches their rows,
and the common vertex bit gives a homeomorphism of the cores.
Write \(H\) for their common homeomorphism type.

A retained row absorbs any bounded tower remainder and any Cantor part:
it contains a clopen even packet of degree above the height bound, and
this packet absorbs these spaces by Lemma~\ref{lem:fragments} and its
proof. It also absorbs one more odd packet of its own tag by shifting
its infinitely many such packets. Both absorptions extend across the
rest of the core. It remains to compare the finite multiplicities of
the packets in the two remainders. In the remainder the multiplicity of
\(G_d\) equals the count of \(M_d\), since in a tower \(D\) is the whole
space and in Cantor space \(D\) is empty. For even \(d\) the common core count is finite;
subtract it from the equal total counts in \(U,V\). For odd \(d=2e+1\)
with no retained row of tag \(e\), the core count is zero. For all
remaining odd degrees, \(H\) absorbs every finite multiplicity of the
packet. This accounts for every summand, giving \(U\cong V\) without
subtracting two infinite cardinals.

Now take tuples in \(\Clop(C_S)\) agreeing on all finite-level formulas.
Their corresponding sign cells agree on every unary finite-level test
by substitution, so the preceding comparison gives homeomorphisms on
all corresponding cells. Gluing proves that the tuples are automorphic.
Apply Lemma~\ref{lem:orbits}.
\end{proof}

One additional consequence of the matching argument will be used on a
branch. It permits exchanging two rows, even when their tags differ.

\begin{lemma}\label{lem:exchange}
If \(U\subseteq C_S\) is clopen and contains the vertex, \(a,b\in S\),
and \(n\in\N\), then \(U\sqcup R_{a,n}\cong U\sqcup R_{b,n}\).
\end{lemma}
\begin{proof}
Use the decomposition in Lemma~\ref{lem:core}. The core compactifies its
retained rows, with infinitely many of every tag in \(S\): a neighborhood
of the vertex contains all but finitely many whole rows. A separate
compact row can be included as one additional index before compactifying,
since a finite extra compact component does not change the neighborhood
basis at the outer vertex. On the two sides the enlarged row sets still
have the same cardinality in each tag. The extra rows have equal cutoff,
so each column has the same finite size on the two sides. Match the
enlarged arrays by Lemma~\ref{lem:match} and restore the identical
remainder by the identity.
\end{proof}
\section{A Boolean sentence associated with a tree}\label{sec:guard}
Write \(\mathcal V=\N^{<\omega}\) for all finite sequences of natural
numbers, \(o\) for the empty sequence and \(vi\) for \(v\) followed by
\(i\). A tree \(T\subseteq\mathcal V\) is closed under initial segments.
An infinite branch is a sequence all of whose finite prefixes lie in \(T\).
We use \emph{all} of \(\mathcal V\) as vertex indices, even for the empty
tree. The tree determines which edges are retained: the edge from \(v\) to
\(vi\) is retained exactly when \(vi\in T\).

Assign distinct natural numbers to the symbols \(p_v\) and \(s_{v,i}\)
using a fixed injection of their disjoint union into \(\N\).
The former are private tags and the latter are edge tags. Put
\[
 S_v=\{p_v\}\ \cup\
 \{s_{u,i}:ui\in T,\ v=u\text{ or }v=ui\},\qquad
 R=R_{p_o,0},\qquad B=\Clop(R).
\]
Only \(S_v\) depends on \(T\). For example, if \(T\) consists of \(o\)
and its one child \((3)\), then \(S_o=\{p_o,s_{o,3}\}\) and
\(S_{(3)}=\{p_{(3)},s_{o,3}\}\). Every other chart has only its private
tag. Extending the tree by a child of \((3)\) gives that chart one more
tag, shared with the new child. In particular \(B\) is a fixed countably
infinite Boolean algebra. A retained edge shares one tag between two charts.
A private tag identifies a single vertex even in the presence of symmetries
of the tree.

We say that a Stone space \(X\) satisfies the following \emph{tree
conditions} when all six requirements hold. The points in the conditions
are determined by the marker sets of \(X\).
\begin{enumerate}
\item For each \(v\), \(A_{p_v}(X)\) is a singleton, written \(\{z_v\}\).
\item The \(z_v\) are pairwise distinct.
\item For \(vi\in T\), \(A_{s_{v,i}}(X)\subseteq\{z_v,z_{vi}\}\).
\item For each \(v\) there is a clopen \(D_v\ni z_v\) homeomorphic to
\(C_{S_v}\).
\item Every clopen \(U\) contains either finitely many or cofinitely many
of the \(z_v\).
\item For each \(j\), every clopen containing infinitely many \(z_v\)
contains a clopen copy of \(G_{2j}\).
\end{enumerate}
By locality and Lemma~\ref{lem:census}, any homeomorphism onto \(D_v\)
takes the vertex to \(z_v\), and \(D_v\) contains no other \(z_w\).
Indeed \(p_w\notin S_v\) for \(w\ne v\). The conditions permit other
points and clopen subspaces outside these charts.

\begin{proposition}\label{prop:guard}
For each \(T\) there is a pure Boolean \(\Pi_{\omega+4}\) sentence
\(\Psi_T\) which, on countable Boolean algebras, expresses the tree
conditions on the Stone space. It has a countable Boolean model, and every
such model is infinite.
\end{proposition}
\begin{proof}
For the sentence, express the first three conditions using
Lemma~\ref{lem:hit}. For example, a closed locus is a singleton if the
whole space meets it and no two disjoint clopens both meet it. Two
singletons differ if a clopen separates them. Inclusion of one closed
locus in the union of two others means every clopen meeting the first meets
at least one of the other two. Each individual test has finite level. For
a clopen \(U\) put
\[
 \mathrm{Fin}(U)=
 \bigvee_{F\subseteq_{\rm fin}\mathcal V}
                 \ \bigwedge_{v\notin F}\neg\Hit_{A_{p_v}}(U).
\]
The inner tests have possibly unbounded finite levels, so this formula is
\(\Sigma_{\omega+1}\). The fifth condition is
\(\forall U\,(\mathrm{Fin}(U)\vee\mathrm{Fin}(U^c))\). For each \(j\)
the sixth says that for all \(U\), either \(\mathrm{Fin}(U)\) holds or
some \(V\le U\) has interval algebra isomorphic to \(\Clop(G_{2j})\).
The latter alternative has finite level by the packet Scott sentence and
relativization to \(V\). Both universal requirements are therefore
\(\Pi_{\omega+2}\).

For the fourth condition use the \(\Pi_{\omega+2}\) chart Scott sentence
of Proposition~\ref{prop:chart} on the interval below \(D_v\), conjoin
\(\Hit_{A_{p_v}}(D_v)\), then quantify existentially over \(D_v\).
For fixed \(v\) this is \(\Sigma_{\omega+3}\); the conjunction over \(v\)
is \(\Pi_{\omega+4}\). The other conditions fit inside the same bound.
Every quantifier used here ranges over elements of the Boolean algebra.
Stone duality gives the claimed equivalences since interval algebras of
countable Boolean algebras are countable.

For an explicit witness take the one-point compactification of the sum
of one chart \(C_{S_v}\) for each \(v\) and, separately, one copy of
\(G_{2j}\) for each \((j,k)\in\N^2\). Its algebra of clopens is countable
by the finite-exception description of a clopen at each compactification.
In each chart the markers are as in Lemma~\ref{lem:census}. No reserve
packet has an odd marker. The outer point lies in every finite derivative
of the closure of isolated points (already the chart vertices do), so it
has no finite degree. For a fixed odd degree all its markers lie in the
at most two charts supporting its tag. This finite union is clopen. Taking
derivatives therefore creates no new anchor at the outer point. The
private and edge requirements follow. A clopen containing the outer point
contains all but finitely many whole components, hence infinitely many
vertices and a copy of each reserve degree. Its complement contains only
finitely many vertices. This gives all six conditions. Finally, any space
satisfying them has infinitely many distinct vertices; a finite Boolean
algebra has a finite Stone space.
\end{proof}

There is one further consequence in an arbitrary space satisfying the
conditions. The set of named vertices has exactly one accumulation point
\(t\). There is at least one by compactness. Two would admit disjoint
clopen neighborhoods each containing infinitely many vertices, contradicting
the finite-or-cofinite condition. Every neighborhood of \(t\) contains
all but finitely many vertices: otherwise its compact clopen complement,
after shrinking the neighborhood to a clopen, would contain another
accumulation point. Each \(D_v\) meets the vertex set in one point, so
\(t\notin D_v\). The point \(t\) controls continuity of the infinite
construction below.
\section{Absorption in every model of the tree conditions}\label{sec:absorb}
Fix \(T\) and a space \(X\) with countably many clopens satisfying the six
conditions. We prove the criterion in both directions. The necessity
argument uses only differences on which the counts are finite. The
sufficiency argument constructs the homeomorphism on clopens small enough
to glue at \(t\).

\subsection{How an additional row forces a branch}
For a clopen \(U\) write \(E_j(U)=\card{M_{2j}(X)\cap U}\) and
\(N_e(U)=\card{P_e(X)\cap U}\); the ambient space will be clear from the
clopen. These quantities are invariant under homeomorphisms of clopens by
locality. Individual \(N_e\)'s can be infinite on a neighborhood of a
vertex.
\begin{lemma}\label{lem:balance}
If \(U\) is a clopen in a row, or a clopen in a chart missing its vertex,
then the \(N_e(U)\) are finite, have finite support in \(e\), and
\[
                  E_j(U)=\sum_e N_e(U)\quad\text{for all sufficiently large } j.
\]
\end{lemma}
\begin{proof}
In a row the only possible root is its own. A clopen missing it meets
finitely many packets, hence misses every sufficiently large even degree.
A clopen containing the root includes every packet except possibly
finitely many, so contains the unique even marker in every sufficiently
large column. This gives the formula in a row. A chart clopen missing its
vertex meets only finitely many rows. Apply the row calculation to its
intersection with each row, take the maximum of the finitely many
thresholds, and add.
\end{proof}

Suppose \(h:X\sqcup R\longrightarrow X\) is a homeomorphism. The added
row has no private anchor; locality shows that \(h\) sends the copy of
each \(z_v\) to \(z_v\). Enumerate \(\mathcal V\) and make the charts
disjoint by taking \(U_v\) equal to \(D_v\) with all earlier \(D_w\)'s
removed. These clopens still contain their respective vertices. In the
source of \(h\), take \(U_v\) for \(v\ne o\), and \(U_o\sqcup R\) for
\(v=o\); call their images \(V_v\). The \(V_v\)'s are also disjoint and
contain their respective \(z_v\)'s. We compare the two families in \(X\);
their unions need not cover \(X\).

Both \(U_v\setminus V_v\) and \(V_v\setminus U_v\) satisfy
Lemma~\ref{lem:balance}, in the second case after pulling back by \(h\)
and allowing the separate clopen piece in \(R\) at \(v=o\). Define
\[
 d_v(e)=N_e(V_v\setminus U_v)-N_e(U_v\setminus V_v)\in\mathbb Z .
\]
This function has finite support contained in \(S_v\). Indeed the two
pieces come from a part of \(D_v\) and, possibly, a part of the added row
of tag \(p_o\in S_o\). Each \(E_j(U_v)\) is finite by
Lemma~\ref{lem:census}. The homeomorphism gives
\(E_j(V_v)=E_j(U_v)+\mathbf1_{\{v=o\}}\), since \(R\) has one marker
in every even column. Cancel the common intersection in this finite
equality and take \(j\) beyond the two thresholds of
Lemma~\ref{lem:balance}. We obtain the first of the following identities:
\begin{align}
 \sum_e d_v(e)&=\mathbf1_{\{v=o\}},\label{eq:sum}\\
 d_v(p_v)&=0,\label{eq:private}\\
 d_v(s_{v,i})+d_{vi}(s_{v,i})&=0 \qquad(vi\in T).\label{eq:edge}
\end{align}

For~\eqref{eq:private}, put \(e=p_v\). Since \(P_e'\subseteq A_e=\{z_v\}\),
the set \(P_e\) has finite intersection with the compact exterior of
\(U_v\), and likewise for \(V_v\). An infinite intersection would have an
accumulation point there. The exteriors have the same number of \(e\)-roots:
\(h\) identifies the latter exterior with the source exterior, whose extra
row is either inside the source sector (\(v=o\)) or has the wrong tag.
On the set \(P_e\), the difference of the indicators of \(V_v,U_v\)
is also the difference of the indicators of their complements in reverse
order. Summing this identity gives zero, exactly~\eqref{eq:private}.
For~\eqref{eq:edge} use \(e=s_{v,i}\), the unions
\(U_v\cup U_{vi}\) and \(V_v\cup V_{vi}\), and
\(P_e'\subseteq\{z_v,z_{vi}\}\). Both exterior counts are finite and equal
by the same argument; \(R\) has no root of tag \(e\). Disjointness
within each pair makes the difference of the union indicators the sum
of the two vertex differences. Only finitely many points contribute to
each difference. This proves~\eqref{eq:edge} without subtracting two
infinite counts.

By~\eqref{eq:sum} the root has a positive coordinate. Its private
coordinate vanishes, so a retained outgoing edge has positive value there.
At its child the same coordinate is negative by~\eqref{eq:edge}. The
sum at the child is zero and has finite support, so another coordinate
is positive. The only possible such coordinate is a retained outgoing
edge: the private coordinate is zero and the unique incoming edge has
already been counted. Induction gives an infinite chain of retained edges
from \(o\). Their child vertices lie in \(T\), and closure under initial
segments puts \(o\) in \(T\) as well. This gives a branch, even when nodes
can have infinitely many children.

\subsection{How a branch absorbs the row}
Conversely let \(v_0=o,v_1,\ldots\) be the successive prefixes of a branch,
where \(v_{n+1}=v_ni_n\). Put \(a_0=p_o\), \(b_n=s_{v_n,i_n}\),
and \(a_{n+1}=b_n\). Both \(a_n\) and \(b_n\) lie in \(S_{v_n}\).

We first describe what can be placed in any clopen neighborhood \(H\) of
both \(z_{v_n}\) and \(t\). In the clopen \(H\cap D_{v_n}\) all but
finitely many whole chart rows remain. Choose a whole row of tag \(b_n\)
there, of cutoff \(k\), and remove it from \(H\cap D_{v_n}\). The
remaining clopen \(L_n\) still contains the vertex. In \(H\setminus
D_{v_n}\), a clopen neighborhood of \(t\), choose successively disjoint
clopen copies of \(G_{2j}\), one for each \(j<k\). This is possible by
condition 6: each chosen packet has no private anchor (its odd markers have
empty second derivative), hence no vertex, and also misses \(t\). After
each choice the remaining neighborhood of \(t\) still has infinitely many
vertices. The removed row together with these packets forms a clopen
\(Q_n\cong R_{b_n,0}\). Indeed the finite added packet components fill the
missing even columns below \(k\), and adjoining finitely many compact
components does not change the neighborhood basis at the row root.
Set \(W_n=L_n\sqcup Q_n\). It misses \(t\), contains just the named vertex
\(z_{v_n}\), and lies in \(H\). By Lemma~\ref{lem:exchange},
\begin{equation}\label{eq:block}
 L_n\sqcup R_{a_n,0}\ \cong\ L_n\sqcup R_{b_n,0}\ \cong\ W_n.
\end{equation}

We must choose the neighborhoods so that these local identifications fit
together. Enumerate a clopen basis \(H_0,H_1,\ldots\) of \(X\). At step
\(n\) take the union of two clopen neighborhoods: the intersection of the
\(H_i\), \(i\le n\), which contain \(t\), and the intersection of those
which contain \(z_{v_n}\) (an empty intersection is \(X\)). Remove the
previously selected \(W_m\), \(m<n\). The result still contains both
points, so the preceding construction applies. The \(W_n\) are disjoint.
For every neighborhood of \(t\), choose \(H_i\) inside it containing
\(t\). Eventually \(z_{v_n}\in H_i\), since the branch has distinct
vertices and \(H_i\) contains cofinitely many of them. For all sufficiently
large \(n\ge i\) both intersections at step \(n\) are contained in
\(H_i\). Thus every neighborhood of \(t\) contains all sufficiently late
whole \(W_n\).

On the punctured spaces, map \(R\sqcup L_0\) to \(W_0\), then
\(Q_{n-1}\sqcup L_n\) to \(W_n\) for \(n\ge1\), using~\eqref{eq:block}
and \(b_{n-1}=a_n\). Use the identity on the remainder outside all the
sites. These are disjoint open partitions of the punctured spaces: away
from \(t\) the sites are locally finite, since they eventually lie in
any clopen neighborhood of \(t\). The piecewise map and its inverse are
therefore continuous. Each original space is the one-point compactification
of its puncture at \(t\); a homeomorphism on the punctures carries compact
sets to compact sets in both directions and hence extends by \(t\mapsto t\).
We have constructed \(X\sqcup R\cong X\).

\begin{theorem}\label{thm:absorb}
For every countable Boolean algebra \(A\models\Psi_T\),
\[
                       A\cong A\times B
       \quad\Longleftrightarrow\quad T\text{ has an infinite branch}.
\]
\end{theorem}
\begin{proof}
Apply the two directions above to the Stone space of \(A\). Its clopens
give a countable basis, and \(\Clop(X\sqcup R)=\Clop(X)\times B\) by
restriction to the two summands.
\end{proof}
\section{Borel comparison and the choice of a sentence}\label{sec:diagonal}
We now have a criterion for when absorption holds. To find a sentence
all of whose countable models have high rank, we need a second test, for
agreement with the product up to a specified logical complexity. Unlike
isomorphism, this second test is Borel at each fixed countable level.

\subsection{Bounded back-and-forth}
For structures \(M,N\) in the operation signature and tuples
\(\bar u,\bar v\) of the same finite length, define
\(F_\alpha(M,\bar u;N,\bar v)\) recursively. Require agreement on all
equalities between terms at the tuples. In addition, for every
\(\beta<\alpha\) and \(k<\omega\), every extension of either tuple by a
\(k\)-tuple must have a matching extension on the other side such that
the extended tuples satisfy \(F_\beta\). At \(\alpha=0\) only the term
equalities remain. Write \(M\equiv_\alpha N\) for this comparison of the
empty tuples.

\begin{lemma}\label{lem:bf}
Tuples satisfying \(F_\alpha\) agree on all \(\Sigma_\alpha\) and
\(\Pi_\alpha\) formulas. Isomorphic structures satisfy the comparison at
all levels on corresponding tuples. For fixed countable \(\alpha\), on
fixed countable domains, the comparison is Borel whenever each atomic
term equality is Borel in the parameters.
\end{lemma}
\begin{proof}
For transfer use induction on the formation of the level. Atomic
agreement handles the quantifier-free case. If a disjunct
\(\exists\bar y\,\theta\) is true, choose its finite witness and use the
extension condition at the level \(\beta<\alpha\) of \(\theta\).
Induction transfers its truth. The reverse transfer uses extension in the
other direction. Negation gives the universal case. An isomorphism and
its inverse supply matching extensions at every stage by induction.
Finally a Borel set is one built from open sets by countable unions and
complements. The term tests form a countable intersection. At a
countable ordinal the smaller ordinals, all block lengths, and all finite
tuples in either domain each form countable sets. The recursion thus uses
only countable intersections and unions of previously Borel relations.
\end{proof}

\subsection{A nonanalytic family of branchless parameters}
Let \(\Baire=\N^{\N}\) have the product topology of discrete copies of
\(\N\). Its basic open sets prescribe finite prefixes. Recall that in
standard Borel spaces (measurable spaces isomorphic to Borel subsets of
Polish spaces), projections of Borel sets are \emph{analytic}. Every
nonempty analytic subset of \(\Baire\) is the range of a continuous map
from \(\Baire\)~\cite[\S14]{kechris}. A Polish space is a separable,
completely metrizable space.

Fix an injective coding \(\langle u,v\rangle\in\N\) of pairs of finite
sequences. Given \(c,x\in\Baire\) form the tree
\[
 s\in T(c,x)\quad\Longleftrightarrow\quad
 \forall m\le\card{s},\quad c(\langle x\res m,s\res m\rangle)\ne0.
\]
Here \(\res m\) takes the prefix of length \(m\), including the empty
prefix at \(m=0\). Put \(T_c=T(c,c)\) and \(W=\{c:T_c\text{ has no
branch}\}\). Membership of a fixed node in \(T_c\) is clopen in \(c\):
the prefix used in the indices and the finitely many values at those
indices can all be fixed in a neighborhood of each \(c\).

\begin{lemma}\label{lem:diag}
The set \(W\) is not analytic.
\end{lemma}
\begin{proof}
Given a continuous \(f:\Baire\to\Baire\), let \(d\) be the zero-one code
marking all pairs \(\langle f(y)\res n,y\res n\rangle\) for all \(y,n\).
Then \(T(d,x)\) has a branch exactly when \(x\) lies in the range of \(f\).
Indeed, \(f(y)=x\) gives the branch \(y\). Conversely, from a branch
\(y\) choose \(z_n\) with \(z_n\res n=y\res n\) and
\(f(z_n)\res n=x\res n\), using the marked pair of length \(n\).
Then \(z_n\to y\), \(f(z_n)\to x\), and continuity gives \(f(y)=x\).
The all-zero parameter belongs to \(W\), so if \(W\) were analytic, some
continuous \(f\) would have range \(W\). At its own code \(d\) the
equivalence just proved says \(T_d\) has a branch if and only if
\(d\in W\), a contradiction.
\end{proof}

\subsection{Why truth of the varying sentence is jointly Borel}
Encode a structure on \(\N\) by two constants and the value tables for
the unary and the two binary operations. The space \(\mathcal T\) of such
tables, a countable product of countable discrete spaces, is Polish.
Evaluation of a fixed term at a fixed tuple depends locally on finitely
many entries, by induction over the term. Boolean tables form a Borel
set by imposing the Boolean identities on all tuples. For a fixed
infinitary formula, truth at a fixed tuple is Borel by induction: the
quantifiers over \(\N\) and the logical conjunctions and disjunctions are
countable.

\begin{proposition}\label{prop:borel}
The set of pairs
\[
 \mathcal G=\{(c,m)\in\Baire\times\mathcal T:
             m\text{ is Boolean and }m\models\Psi_{T_c}\}
\]
is Borel.
\end{proposition}
\begin{proof}
The singleton, distinctness, finite-or-cofinite and reserve requirements
use fixed formulas. For an edge the condition is the implication from
membership of its child in \(T_c\) to a fixed inclusion test. Intersecting
over edges is Borel. The chart condition is different. The chart to be
recognized depends on \(c\), and Proposition~\ref{prop:chart} alone
supplies no uniform choice of Scott sentence. We check the semantic
condition using a single countable supply of codes.

Use all clopens of the universal chart \(C_{\N}\) as codes. Restrict a
code to its vertex and to the rows \((e,k)\) with \(e\in S_v(T_c)\).
This gives a surjection onto \(\Clop(C_{S_v(T_c)})\): a clopen of the
smaller chart extends by assigning its vertex bit constantly on all
discarded rows. The finite-exception criterion proves the extension
clopen. Restriction preserves all Boolean operations. For any fixed two
term values on a tuple of codes, their restrictions are equal exactly when
their vertex bits agree and
\[
 \forall e,k,\qquad e\in S_v(T_c)\ \Longrightarrow\
       \text{the two values agree throughout row }(e,k).
\]
The row values here are taken once in the universal chart. Whether
\(e\in S_v(T_c)\) is Borel: the tag is private, unused, or names a
specific edge, whose incidence and retention determine the answer.
This makes every atomic term equality Borel in \(c\).

On the other side, for a proposed chart clopen \(d\) in \(m\), take codes
\(a\in\N\) for \(a\wedge d\) in the interval below \(d\). This too is
surjective. Evaluate relative terms with top \(d\) and complement
\(d\wedge\neg x\); every term test is Borel in the table. On these two
fixed countable sets of codes apply the recursion of Lemma~\ref{lem:bf}
using the just described term equalities. The result computes the true
comparison on the represented algebras. At stage zero this is the
definition of their term diagrams; at each extension step, a tuple of
actual elements can be lifted to codes by surjectivity, and any tuple of
codes represents actual elements. Repeated codes for the same element
cause no change.

At level \(\omega+2\) this comparison is precisely the desired chart
isomorphism test. An isomorphism implies comparison. Conversely comparison
transfers the chart's \(\Pi_{\omega+2}\) Scott sentence to the countable
interval, so the two algebras are isomorphic. Combine with the fixed test
that \(d\) meets the private anchor, take the countable union over \(d\),
then the countable intersection over \(v\). This proves the joint Borel
assertion on Boolean tables.
\end{proof}

\subsection{Excluding every low-rank model at once}
Fix any \(b<\omega_1\). The product \(m\times B\) can be read on the fixed
countable domain \(\N\times B\), using operations coordinate by coordinate.
Its atomic term tests are Borel in \(m\). Form the set
\[
 P_b=\{c:\exists m\in\mathcal T,\ (c,m)\in\mathcal G
                     \ \text{and}\ m\not\equiv_{b+1}m\times B\}.
\]
Proposition~\ref{prop:borel} and Lemma~\ref{lem:bf} make \(P_b\) a
projection of a Borel set, hence analytic. If \(T_c\) has a branch, every
guarded \(m\) absorbs the product by Theorem~\ref{thm:absorb}. An
isomorphism implies every bounded comparison, so \(c\notin P_b\).
Thus \(P_b\subseteq W\). Lemma~\ref{lem:diag} gives a parameter
\(c\in W\setminus P_b\). This choice is the nonconstructive step deciding
which sentence to use for \(b\).

Every countable Boolean \(A\models\Psi_{T_c}\) is infinite by
Proposition~\ref{prop:guard}, hence isomorphic to a table on \(\N\).
The choice \(c\notin P_b\) gives \(A\equiv_{b+1}A\times B\). The choice
\(c\in W\) gives \(A\not\cong A\times B\). Consequently every
\(\Pi_{b+1}\) sentence true of \(A\) is true of a nonisomorphic countable
Boolean algebra, namely \(A\times B\). The product need not itself
satisfy \(\Psi_{T_c}\): the Scott sentence must distinguish \(A\) from
every nonisomorphic countable structure. This explains the role of
projecting the existence of even one failed comparison.

\begin{proof}[Completion of the proof of Theorem~\ref{thm:main}]
Conjoin \(\Psi_{T_c}\) with the universal Boolean identities. Explicitly
one may use commutativity and associativity of both lattice operations,
both absorption identities, distributivity of \(\vee\) over \(\wedge\),
and \(0\vee x=x,\ x\vee1=1,\ x\wedge\neg x=0,\ x\vee\neg x=1\).
The lattice laws give a lattice under \(x\le y\) iff \(x\wedge y=x\);
the remaining laws give its bounds, distributivity and Boolean complement.
Their universal closures form a \(\Pi_1\) sentence, so the conjunction
\(\varphi_b\) remains \(\Pi_{\omega+4}\). The Boolean witness of
Proposition~\ref{prop:guard} satisfies it. Any operation structure
satisfying the conjunction is a Boolean algebra. If its zero equaled its
one it would be a singleton, hence a countable model of the tree
conditions, contradicting Proposition~\ref{prop:guard}. This also proves
nontriviality for unrestricted model cardinalities. Finally the product
comparison just proved supplies a nonisomorphic countable competitor
even within Boolean algebras. Order and operation isomorphisms coincide:
the operations are determined by the order, and the order is determined
by meet. Thus a Scott sentence isolating any countable model in the raw
operation signature cannot have level \(\Pi_{b+1}\).
\end{proof}
\clearpage
\section*{Guide to the formal proof}
\addcontentsline{toc}{section}{Guide to the formal proof}
The links below point to the Lean source for the corresponding steps.
The endpoint theorem is named
\texttt{BooleanScott.pure\_boolean\_counterexample}. Its quantifiers run over
arbitrary operation structures for the implication to the Boolean laws
and over countable operation structures for the failure of Scott isolation.

\begin{center}
\begin{tabularx}{\linewidth}{@{}lX@{}}
\toprule
Argument in this paper & Entry points in the formalization\\
\midrule
Packet types (\S\ref{sec:packets}) &
\code{MixedPacket}{Mixed packets}, \code{PacketClassification}{classification of clopens},
\code{PacketScott}{finite-level isolation}\\
Charts (\S\ref{sec:charts}) &
\code{ArrayHomeomorphism}{matching the columns},
\code{ChartCoreMatching}{matching the cores},
\code{ChartIsolation}{chart Scott sentences},
\code{RowExchange}{row exchange}\\
Tree conditions (\S\ref{sec:guard}) &
\code{TreeWitness}{the explicit witness},
\code{GuardGeometry}{geometric conditions},
\code{GuardSentence}{compilation to a Boolean sentence}\\
Absorption (\S\ref{sec:absorb}) &
\code{AbsorptionNecessity}{extracting a branch},
\code{BranchBlocks}{allocating the clopen blocks},
\code{BranchSchedule}{choosing shrinking sites},
\code{BranchGluing}{gluing along the branch}\\
Choice of parameter (\S\ref{sec:diagonal}) &
\code{Diagonal}{diagonal family}, \code{ChartPresentation}{chart codes},
\code{GuardBorel}{joint Borel relation}, \code{Reduction}{rank obstruction},
\code{PureCounterexample}{theorem in the pure signature}\\
\bottomrule
\end{tabularx}
\end{center}

\paragraph{Contributions and review.}
OpenAI Codex developed the mathematical argument and Lean formalization under
the direction of Serafim Batzoglou. Batzoglou directed the proof search and
requested adversarial checks. The formalization passed the documented Lean
verification. Batzoglou has not yet read the proof in full or independently
verified that the formal theorem resolves the original question. Independent
expert review is pending.
This manuscript is licensed under
\href{https://creativecommons.org/licenses/by/4.0/}{Creative Commons Attribution 4.0 International}.

\begin{thebibliography}{9}
\addcontentsline{toc}{section}{References}
\bibitem{ght} D. Gonzalez and M. Harrison-Trainor.
\href{https://arxiv.org/html/2411.12084v1}{Scott spectral gaps are bounded for linear orderings}.
arXiv:2411.12084v1 (2024). Theorem~1.5 and Question~1.9.
\bibitem{ht} M. Harrison-Trainor.
\href{https://doi.org/10.1016/j.aim.2018.03.012}{Scott ranks of models of a theory}.
\emph{Advances in Mathematics} \textbf{330} (2018), 109--147.
\bibitem{kechris} A. S. Kechris.
\href{https://doi.org/10.1007/978-1-4612-4190-4}{\emph{Classical Descriptive Set Theory}}.
Graduate Texts in Mathematics 156, Springer, 1995.
\bibitem{ketonen} J. Ketonen.
\href{https://annals.math.princeton.edu/1978/108-1/p02}{The structure of countable Boolean algebras}.
\emph{Annals of Mathematics} (2) \textbf{108} (1978), 41--89.
\bibitem{montalban} A. Montalb\'an.
\href{https://doi.org/10.1090/proc/12669}{A robuster Scott rank}.
\emph{Proceedings of the American Mathematical Society} \textbf{143} (2015), 5427--5436.
\bibitem{scott} D. Scott.
\href{https://doi.org/10.1016/B978-0-7204-2233-7.50036-8}{Logic with denumerably long formulas and finite strings of quantifiers}.
In \emph{The Theory of Models}, North-Holland, 1965, 329--341.
\bibitem{stone} M. H. Stone.
\href{https://doi.org/10.1090/S0002-9947-1936-1501865-8}{The theory of representations for Boolean algebras}.
\emph{Transactions of the American Mathematical Society} \textbf{40} (1936), 37--111.
\end{thebibliography}
\end{document}
